\documentclass[10pt]{article}

\usepackage[utf8]{inputenc}

\usepackage[T1]{fontenc}

\usepackage{amsmath}

\usepackage{amsfonts}

\usepackage{amssymb}

\usepackage[version=4]{mhchem}

\usepackage{stmaryrd}

\usepackage{bbold}



\begin{document}

поверхностей. Этот класс метрич. пространства можно получить, присоединяя к двумерным римановым пространствам двумерные метризованные многообразия, метрика к-рых в окрестности каждой точки допускает равномерное приближение римановыми метриками с ограниченными в совокупности интегралами от абсолютной гауссовой кривизны.



П р о с т р а н с в а к р и в иныне боль ш $K$ - теория полных метрич. многообразий с внутренней метрикой, в к-рых сумма верхних углов треугольников, составленных из кратчайших, не превосходит суммы углов треугольника на плоскости постоянной кривизны $K$ с такими же длинами сторон (кроме того, предполагается, что любые две точки можно соединить единственной кратчайшей).



См. также Геодезических геометрия, Конформная геометрия, Риманово пространство обобщенное.



Лит.. [1] Р и м а н Б., Соч., пер. с нем., М.-Л., 1948; [2] Р а ші е в с к и й П. К. Риманова геометрия и тензорный анализ, 3 изд., М., 1967; [З̉] У й 3 е н х а $\mathbf{p}$ т Л. П., Риманова геометрия, пер. с англ., М., 1948; [4] Г р о м о л Д., К л и н г е нбе р г В., М е й е р В., Риманова геометрия в целом, пер. с нем., М., 1971; [5] А'ле к с а н д р о в А. Д., Внутренняя геометрия выпуклых поверхностей, М.-Л., 1948; [6] Б у р аг о I0. Д., 3 а л г а л л е р В. А., "Успехи матем. наук", 1977, т. 32, в. 3, с. 3-55; [7] М и л н о р Д ж., Теория Морса, пер. с англ., М., 1965 ; [8] К а $\mathbf{~ т ~ а ~ н ~ Э . , ~ Г е о м е т р и я ~ р и м а н о в ы х ~}$ пространств, пер. с франц., м.-J., 1936 ;'[9] K u l k a r n i R. S., "Ann. Math.", 1970, v. 91, № 2, p. 311 -31; [10] W o l f J. A., Spaces of constant curvature, N. Y., 1967. B. A. Топологов.



РИМАНОВА ГЕОМЕТРИЯ В ЦЕЛОМ - раздел римановой геометрии, изучающий связи между локальными и глобальными характеристиками римановых многообразий (р. м.). Термин Р. г. в ц. обычно относят к определенному кругу проблем и методов, характерных для геометрии в целом. Основное место в Р. г. в ц. занимает изучение связи между кривизной и топологией р. м. При әтом исследуются вопрос о топологическом и метрич. строении р. м. с данными условиями на кривизну и вопрос о существовании на заданном гладком многообразии римановой метрики с предписанными свойствами кривизны (секционной кривизны $K_{\sigma}, Р$ чччи кривизны Ric, скалярной кривизны $K_{\mathrm{ck}}$ ). Бо́льшая qасть полученных результатов относится к пространствам с кривизнами постоянного знака. Р. г. в ц. тесно соприкасается с теорией однородных пространсте и вариационной теорией геодезических линий. О подмногообразиях Р. м. см. в статьях Изометрическое погружение и Погруженных многообразий геометрия.



Методы Р. г. в ц. носят синтетич. характер. Наряду с локальной дифференциальной геометрией широко используются теория дифференциальных уравнений и Морса теория. Но основные достижения связаны с нахождением удачных конструкций, напр. построением замкнутых геодезических, минимальных пленок или пленок из геодезических, орисфер, выпуклых множеств. Исследованию топологии р. м. обычно предшествует изучение их метрич. свойств. Последнее часто осуществляется путем сравнения р. м. с подходящим әталонным пространством (см. ниже теоремы сравнения).



Т о поло г и п с к о е с т р о в и е. Для замкнутых поверхностей связь между кривизной и топологией по существу исqерпывается формулой Гаусса Бонне. Среди замкнутых поверхностей метрику положительной кривизны могут нести только сфера $S^{2}$ и проективная плоскость $P^{2}$, а нулевой кривизны - тор и бутылка Клейна. Строение р. м. размерности $n>2$ известно хуже (1983). Вот примеры известных теорем.



Полное односвязное р. м. $M^{n}$ с $K_{\sigma} \leqslant 0$ диффеоморфно $\mathbb{R}^{n}$ (т е о ре м а А д а м а р а - К а р т а н a), приqем для любой точки $x \in M^{n}$ экспоненциальное отображение $\exp _{x}$ есть дифффеоморфизм касательного пространства $T_{x} M^{n}$ на $M^{n}$.



Для замкнутых р. м. с $K_{\sigma}>0$ имеет место следующая т е о р е м о с ф е р е: полное р. м. $M^{n}$ с $0<\delta \leqslant$ $\leqslant K_{\sigma} \leqslant 1$ наз. $\delta-3$ а пा е м л е н н ы м; если оно односвязно и $\delta>1 / 4$, то $M^{n}$ гомеоморфно $S^{n}$. Для четных $n$ граница здесь точная: при $\delta=1 / 4$ существуют $M^{n}$, негомеоморфные $S^{n}$, это - симметрические пространства ранга 1 и только они. Для нечетных $n$ теорема о гомеоморфности $M^{n}$ и $S^{n}$ верна и при $\delta=1 / 4$. При $n<7, n \neq 4$, гомеоморфность сфере $S^{n}$ влечет диффеоморфность. При $n \geqslant 7$ диффеоморфность сфере установлена при более жестком, чем в теореме о сфере, защемлении (достаточно взять $\delta>0,87$, а при $n \rightarrow \infty$ взять $\delta>0,66)$. Известно также, что при ещце более жестком зацемлении (достаточно $\delta>0,98$ и $\delta>0,66$ при $n \rightarrow \infty$ ) неодносвязное $M^{n}$ дифффеоморфно пространству постоянной кривизны (факторпространству $S^{n}$ по дискретной подгруппе изометрий). Есть ряд результатов об условиях на $K_{\sigma}$, обеспечивающцих гомеоморфизм симметрич. пространству ранга 1 (cM. [14], [15]).



Открытое, т. е. полное некомпактное, р. м. с $K_{\sigma}>0$ всегда диффеоморфно $\mathbb{R}^{n}$. Множество $X \subset M^{n}$ наз. а б с о л ю т о в ы п у л ы м, если каждая геодезическая с концами в $X$ вся лежит в $X$. Пусть $M^{n}-$ открытое р. м. с $K_{\sigma} \geqslant 0$, тогда в $M^{n}$ суцествует вполне геодезическое абсолютно выпуклое замкнутое подмногообразие $N$ такое, что $M^{n}$ диффеоморфно пространству $v(N)$ нормального расслоения $N$ в $M^{n}$ (если $K_{\sigma}>0$, то $\operatorname{dim} N=0)$. В случаях $\operatorname{dim} N=1$ или $n-1$, а для однородных пространств - всегда, $M^{n}$ даже изометрично $v(N)$ со стандартной метрикой нормального расслоения. При $n \leqslant 3$ это дает полную классификацию открытых p. м. с $K_{\sigma} \geqslant 0$.



П р я м о й наз. полная геодезическая, кратчайшая на любом своем участке. Т е о р е м а ц ц л и н д р е: открытое $M^{n} \mathrm{c} K_{\sigma} \geqslant 0$ изометрично прямому метрич. произведению $M^{n-k} \times \mathbb{R} k, 0 \leqslant k \leqslant n$, где $M^{n-k}$ не содержит прямых. Условие $K_{\sigma} \geqslant 0$ здесь можно заменить на Ric $\geqslant 0$



Фунд а мента льна я г р уп па. При $K_{\sigma}>0$ и qетном $n$ замкнутое $M^{n}$ либо ориентируемо и односвязно, либо неориентируемо и фундаментальная группа $\pi_{1}\left(M^{n}\right)=\mathbb{Z}_{2} ;$ при неqетном $n$ оно всегда ориентируемо, но о $\pi_{1}\left(M^{n}\right)$ за пределами теоремы о сфере мало что известно. Даже для $M^{n}$ постоянной кривизны $K_{\sigma}=1$ полное описание возможных строений $\pi_{1}\left(M^{n}\right)$ для нечетных $n$ оказалось трудной задачей (см. [9]).



Если $K_{\sigma}=0$, то универсальное накрывающее $\tilde{M}^{n}$ пространства $M^{n}$ изометрично $\mathbb{R}^{n}$ и фундаментальная группа $\pi_{1}\left(M^{n}\right)$ изоморфна дискретной группе изометрий $\mathbb{R}^{n}$ без неподвижных точек; она содержит подгруппу сдвигов конечного индекса. (Тем самым $M^{n}$ допускает конеqное изометрич. накрытие плоским тором.)



Если в $\tilde{M}^{n}$ все $K_{\sigma} \leqslant 0$, то $\tilde{M}^{n}$ диффеоморфно $\mathbb{R}^{n}$. Поәтому все гомотопич. группы $\pi_{i}\left(M^{n}\right)$ для $i>1$ тривиальны, и гомотопич. тип определяется $\pi_{1}\left(M^{n}\right)$. Если $K_{\sigma}<0$, то $\pi_{1}\left(M^{n}\right)$ полностью некоммутативна в том смысле, что любая абелева (и даже любая разрешимая) ее подгруппа является бесконеqной циклической. В случае $K_{\sigma} \leqslant 0$ известно следующее. Пусть $\Gamma$ - разрешимая подгруппа $\pi_{1}\left(M^{n}\right)$. Тогда $\Gamma$ изоморфна дискретной группе изометрий $\mathbb{R}^{n}$ (без неподвижных точек) и $M^{n}$ содержит компактное вполне геодезич. подмногообразие, изометричное $\mathbb{R} n / \Gamma$. Вместо $K_{\sigma} \leqslant 0$ достаточно при этом потребовать отсутствия на геодезических сопряженных точек.



Для двух многообразий одинаковой постоянной отрицательной кривизны и одинаковой размерности $n \geqslant 3$ изоморфизм $\pi_{1}$ влеqет изометрию (т е о р е м а $\mathrm{M}$ о сT 0 B a).



Римановы многообразия, для к-рых $\max \left|K_{\sigma}\right| \times$ $\times \operatorname{diam} M^{n}<\varepsilon$, наз. $\varepsilon-\Pi$ л о с к и м и. Такие многообразия могут при произвольном $\varepsilon>0$ топологически отлиqаться от локально плоских. Для них при любом $n \geqslant 2$ существует такое $\varepsilon_{n}$, что для $\varepsilon_{n}$-плоского $M^{n}$ в $\pi_{1}\left(M^{n}\right)$





\end{document}